\par\noindent\textbf{(1)} 小应变 \(\varepsilon_{ij}=(\partial_iu_j+\partial_ju_i)/2\)，应力 \(\sigma_{ij}=\lambda\delta_{ij}\nabla\cdot\boldsymbol u+2\mu\varepsilon_{ij}\)。动量守恒 \(\rho u_{i,tt}=\partial_j\sigma_{ij}\) 得Navier方程 \(\rho\boldsymbol u_{tt}=(\lambda+\mu)\nabla(\nabla\cdot\boldsymbol u)+\mu\nabla^2\boldsymbol u\)。
\par\noindent\textbf{(2)} 把位移写成无旋 \(\boldsymbol u_L\) 与无散 \(\boldsymbol u_T\) 部分。平面纵波 \(\boldsymbol u\parallel\boldsymbol k\) 给 \(\omega^2=(\lambda+2\mu)k^2/\rho\)，所以 \(c_L=\sqrt{(\lambda+2\mu)/\rho}\)；横波 \(\boldsymbol u\perp\boldsymbol k\) 给 \(c_T=\sqrt{\mu/\rho}\)。
\par\noindent\textbf{(3)} 单位体积弹性势能 \(W=\lambda(\operatorname{tr}\varepsilon)^2/2+\mu\varepsilon_{ij}\varepsilon_{ij}\)。将应变分解为体积与无迹部分，得 \(W=(K/2)(\operatorname{tr}\varepsilon)^2+\mu\varepsilon_{ij}^{\rm dev}\varepsilon_{ij}^{\rm dev}\)。对任意应变正定要求 \(K>0,\mu>0\)，从而 \(c_L>c_T>0\)。
\par\noindent\textbf{(4)} 近不可压极限 \(K\gg\mu\) 时 \(c_L\sim\sqrt{K/\rho}\) 很大而 \(c_T\) 仍有限；表面存在由纵横体波耦合形成的Rayleigh表面波，其速度低于 \(c_T\)，不可把表面波直接等同任一体波。
